\section{Introduction}

% \renewcommand{\thefootnote}{\fnsymbol{footnote}}
% \footnotetext[1]{The corresponding authors are Lu Chen and Kai Yu.}

The primary goal of tool learning is to enable Large Language Models (LLMs; \citealp{geminiteam2023gemini,achiam2023gpt,dubey2024llama}) to understand and effectively use external tools~\cite{qin2024toollearningfoundationmodels}. By integrating diverse tools, LLMs can enhance their ability to tackle complex natural language processing (NLP) tasks~\cite{schick2023toolformer,hao2023toolkengpt, hsieh2023tool, tang2023toolalpaca}. Furthermore, recent studies~\citep{qin2023webcpm, NEURIPS2022_82ad13ec_webshop, cai2024large} have explored how LLMs can interact with real-world tasks through tool use, establishing tool learning as a critical bridge between LLMs and the physical world.

However, a significant challenge arises when LLMs are required to reliably invoke external tools. A major issue is \textbf{Tool Hallucinations}~\citep{patil2023gorilla}, where the model either selects inappropriate tools or misuses them. 
% This paper introduces the concept of \textbf{reliable tool calling}, systematically analyzes the causes and types of tool hallucinations, and proposes strategies to mitigate these errors, thereby enhancing the reliability of LLM tool interactions.
Addressing tool hallucinations is crucial due to their significant impact on both task performance and system reliability. First, hallucinations during task execution can directly affect real-world systems, as tools act as the interface between LLMs and the physical world. This is particularly critical in areas such as database operations, robotic control, and scientific experimentation, where incorrect tool use can cause tangible damage and pose greater risks than textual hallucinations~\citep{ji2023survey}. Second, tool hallucinations produce \textit{task hallucination} that are harder to detect. Since hallucinations occur during task execution, they often result in misleading final outputs, undermining user trust in the system’s reliability. Finally, from an operational standpoint, hallucinated tool calls increase computational costs, and reduce efficiency by generating redundant or faulty tool invocations. These challenges highlight the need for a systematic approach to understanding and mitigating tool hallucinations.

To address this challenge, we pioneer a comprehensive evaluation framework for \textbf{Tool Reliability}, which consists of three key components. First, we systematically define and categorize tool hallucination, introducing an automated evaluation framework and corresponding metrics. Tool hallucination is classified into two main types and four subtypes, as illustrated in Figure~\ref{fig:hallucination_type}. \textit{Tool selection hallucination} occurs when the model incorrectly chooses a tool unrelated to the task or calls it at an inappropriate time. \textit{Tool usage hallucination} involves errors in applying the selected tool, such as providing incorrect parameters or fabricating non-existent ones.  To assess these hallucinations, we design a hybrid evaluation approach that combines rule-based methods with LLM evaluators, enabling automated and accurate detection of hallucination instances.

Second, we introduce a task-level reliability measurement framework that addresses the cascading effects of tool hallucinations on task outcomes. Since tool hallucinations often lead to hallucinated final outputs, we refine the traditional task success rate by incorporating hallucination-aware corrections. Specifically, we propose the \textbf{Re}liable \textbf{P}ass \textbf{R}ate (\textbf{RePR}), which discounts task outcomes that are influenced by hallucinated tool calls. Furthermore, to capture the trade-offs between task success and resource efficiency, we introduce the \textbf{Benefit-cost Utility} metric. This metric not only rewards successful task completion but also penalizes hallucinations and excessive tool-calling costs.

Third, we develop \textbf{RelyToolBench}, a specialized benchmark derived from StableToolBench~\cite{guo2024stabletoolbench}, to rigorously evaluate tool hallucinations. RelyToolBench includes two additional types of subsets: one addressing tool selection hallucinations through tool-task mismatches, and another focusing on tool usage hallucinations through parameter deficiencies. These subsets enable precise assessment of model behavior in tool invocation tasks, providing granular insights into tool hallucination and task reliability.

% To mitigate tool hallucinations, we propose a reliability-focused alignment framework. For the first time, this framework expands the action space to include \textbf{indecisive actions}, such as engaging in dialogue with the user or switching tools. Unlike traditional \textbf{decisive actions} (direct tool calls), which can lead to hallucinations or errors, indecisive actions enable the model to better assess situations and avoid erroneous decisions. 

To mitigate tool hallucinations and enhance task reliability, we propose \textbf{Relign}—a \textbf{Re}liability a\textbf{lign}ment framework for tool calling. Tool invocation can be regarded as a sequential decision-making process in which LLMs execute specific actions. Unlike conventional reinforcement learning tasks, where all actions within the action space are generally valid, tool invocation involves the risk of issuing invalid or hallucinated tool calls. Prior works~\cite{qin2023toolllm, yu2024steptool} primarily focus on improving the model’s ability to make better tool selections from a predefined set of \textbf{decisive actions} (i.e., directly invoking a tool). However, when the preconditions for tool usage are not met—such as when the user fails to provide sufficient parameters or when the selected tool is unsuitable for the given task—any action chosen from the decisive action space inevitably leads to an invalid tool call and results in hallucination. To address this issue, Relign introduces the concept of an \textbf{indecisive action space} within the tool-use environment, allowing the model to proactively learn when to defer tool invocation, engage in clarification dialogues with the user, or switch to an alternative tool. This expansion of the action space enables LLMs to make more informed decisions, reducing hallucinations while improving overall system robustness.

Additionally, Relign incorporates optimization techniques like Supervised Fine-tuning (SFT; \citealp{zhang2023instruction}) and Direct Preference Optimization (DPO; \citealp{rafailov2024direct}), along with a specialized data synthesis algorithm to generate reliability-focused training data. Experiments on RelyToolBench demonstrate that Relign effectively reduces tool hallucinations and significantly improves model reliability. Our contributions are as follows:
\begin{itemize}
\item We systematically define and categorize tool hallucinations into two major types with four subtypes, proposing an automated evaluation framework.
\item We introduce the RePR metric for hallucination-aware task evaluation, develop Benefit-cost Utility for cost-sensitive analysis, and curate RelyToolBench with specialized test cases for rigorous reliability assessment.
\item We propose Relign, the first tool reliability alignment framework, supported by a data synthesis pipeline and preference optimization techniques to reduce tool hallucinations and improve overall reliability.
\end{itemize}